\documentclass[a4paper,10pt]{article}
\usepackage{amsmath,amssymb}
\usepackage{fontspec}
\setmainfont{Liberation Serif}
\usepackage{graphicx}
\usepackage{eso-pic}
\usepackage{tikz}
\usepackage[most]{tcolorbox}
\usepackage{geometry}
\usepackage{fancyhdr}
\usepackage{lastpage}
\usepackage{booktabs}
\usepackage{listings}
\usepackage{enumitem}
\usepackage{microtype}

\geometry{
    a4paper,
    top=1.5cm,
    bottom=1.8cm,
    left=1.2cm,
    right=1.2cm
}

% Persistent Unremovable Canvas Watermark on Every Page
\newsavebox{\watermarkbox}
\sbox{\watermarkbox}{%
    \begin{tabular}{c}
        \includegraphics[width=8cm]{/home/johan/GATE/GATE_CS_PYQ_DB/assets/studiva-butterfly-icon.jpg} \\[-0.5em]
        {\fontsize{18}{22}\selectfont\bfseries\scshape www.studiva.co.in}
    \end{tabular}%
}
\AddToShipoutPictureBG{%
    \begin{tikzpicture}[remember picture, overlay]
        \node[opacity=0.12] at (current page.center) {\usebox{\watermarkbox}};
    \end{tikzpicture}%
}

% Header and Footer
\pagestyle{fancy}
\fancyhf{}
\fancyhead[L]{\small\scshape GATE 2018--2025 CS/IT $\cdot$ \textbf{Engineering Mathematics}}
\fancyhead[R]{\small\scshape Practice Question Bank with Solutions}
\fancyfoot[L]{\scriptsize\textsc{Copyright \copyright\ 2026} Studiva\textsuperscript{TM}}
\fancyfoot[C]{\small Page \thepage\ of \pageref{LastPage}}
\fancyfoot[R]{\scriptsize\texttt{www.studiva.co.in}}
\renewcommand{\headrulewidth}{0.4pt}
\renewcommand{\footrulewidth}{0.4pt}

% Code Listing Style (Monochrome)
\lstset{
    basicstyle=\footnotesize\ttfamily,
    numbers=none,
    frame=none,
    tabsize=2,
    breaklines=true,
    showstringspaces=false
}

\begin{document}

% ==================== PAGE 1: IMPRINT & TITLE PAGE ====================
\thispagestyle{empty}
\vspace*{0.5cm}
\begin{center}
    \includegraphics[width=2.5cm]{/home/johan/GATE/GATE_CS_PYQ_DB/assets/studiva-butterfly-icon.jpg}\\[0.6em]
    {\LARGE \textbf{GATE 2018--2025 COMPUTER SCIENCE \& IT}}\\[0.3em]
    {\Large \textbf{PREVIOUS YEAR QUESTION BANK WITH SOLUTIONS}}\\[0.4em]
    {\large \textsc{\textbf{Subject Wise: Engineering Mathematics}}}\\[0.2em]
    {\normalsize \textsc{Combined 2018--2025 Examination Papers $\cdot$ Comprehensive Edition ($\mathbf{36}$ Questions)}}\\[1.2em]
    \rule{0.6\textwidth}{0.6pt}
\end{center}

\vspace{1em}
\noindent\begin{tabular}{ll}
\textbf{Author / Compiler:} & Rupam Ghosh \\
\textbf{Publisher \& Platform:} & Studiva\textsuperscript{TM} (\texttt{www.studiva.co.in}) \\
\textbf{Official Inquiries:} & \texttt{contact@studiva.co.in} \\
\textbf{Edition:} & First Standard Edition (2026) \\
\textbf{Coverage:} & Complete GATE CS/IT Papers (2018, 2022, 2023, 2024 CS1, 2024 CS2, 2025 CS1, 2025 CS2) \\
\textbf{Syllabus Reference:} & Official GATE 2026 Syllabus (Organizing Institute: IIT Guwahati) \\
\end{tabular}

\vspace{1.5em}
\begin{tcolorbox}[enhanced, standard jigsaw, opacityback=0, colframe=black!40, arc=1mm, boxrule=0.5pt, left=10pt, right=10pt, top=8pt, bottom=8pt]
\textbf{\large COPYRIGHT \& INTELLECTUAL PROPERTY NOTICE}\\[0.6em]
Copyright \copyright\ 2026 Rupam Ghosh \& Studiva\textsuperscript{TM}. All Rights Reserved.\\[0.5em]
This educational compilation, problem organization, taxonomy classification, LaTeX typesetting, hints, and explanatory solutions are the proprietary intellectual property of \textbf{Rupam Ghosh} and \textbf{Studiva\textsuperscript{TM}}.\\[0.6em]
\textbf{Terms of Use \& Attribution Rules:}
\begin{itemize}[leftmargin=*, nosep]
    \item \textbf{Personal Educational Use:} Permitted freely for individual students and educators preparing for competitive examinations.
    \item \textbf{Mandatory Attribution:} Any reproduction, republication, hosting, or sharing (digital or physical) must prominently credit \textbf{Rupam Ghosh \& Studiva\textsuperscript{TM}} and include an active, direct link to \texttt{https://www.studiva.co.in}.
    \item \textbf{Commercial Restrictions:} Commercial resale, sublicensing, course bundling, or inclusion in paid test series without prior written permission is strictly prohibited.
    \item \textbf{Anti-Tampering Notice:} Removing, altering, or obscuring copyright notices, page watermarks, or website links from this document is strictly prohibited.
\end{itemize}
\vspace{0.6em}
\textbf{Fair Use Disclaimer:}\\
The underlying examination questions originate from the Graduate Aptitude Test in Engineering (GATE), conducted by the Indian Institutes of Technology (IITs) and Indian Institute of Science (IISc). Question texts are reproduced for non-commercial educational study under Fair Use doctrines. No copyright is claimed over the original GATE examination question texts.
\end{tcolorbox}

\vspace{1.2em}
\begin{center}
\small\itshape
Practice Recommendation: Use an opaque sheet of paper or ruler to cover the right-hand column during mock attempts. Tick your chosen answers directly in the check boxes on the left, then slide the cover to check the solution and hint.
\end{center}

\newpage
% ==================== PAGE 2 ONWARDS: QUESTION BANK ====================

\vspace{8pt}
\subsection*{Linear Algebra}
\vspace{-4pt}
\noindent
\begin{tcolorbox}[
    blanker,
    sidebyside,
    sidebyside align=top,
    lefthand width=0.55\textwidth,
    top=4pt, bottom=4pt, left=0pt, right=0pt,
    middle=16pt,
    underlay={
        \draw[dashed, black!80, line width=0.8pt, dash pattern=on 4pt off 3pt]
            ([xshift=0.55\textwidth+8pt]frame.north west) --
            ([xshift=0.55\textwidth+8pt]frame.south west);
    }
]
\textbf{Q.1} \quad {\footnotesize\textsc{[NAT $\cdot$ 1 Mark(s) $\cdot$ No Neg]}} \hfill {\footnotesize\textbf{GATE 2018 CS} \quad \tiny\texttt{2018\_CS\_GA\_Q17}}
\vspace{2pt}\hrule height 0.3pt\vspace{4pt}

Consider a matrix $A = u v^T$ where $u = \begin{bmatrix} 1 \\ 2 \end{bmatrix}$ and $v = \begin{bmatrix} 1 \\ 1 \end{bmatrix}$. Note that $v^T$ denotes the transpose of $v$. The largest eigenvalue of $A$ is \_\_\_\_\_\_\_\_.

\par\vspace{6pt}\noindent\textbf{Your Answer:} \fbox{\rule{0pt}{1.3ex}\hspace{3cm}}

\tcblower

\textbf{Official Key:} \textbf{\large 3 to 3} \hfill \textit{\small [Easy]}
\vspace{2pt}\hrule height 0.3pt\vspace{3pt}
\textbf{Topic:} Linear Algebra \\
\textbf{Subtopic:} \textit{Eigenvalues and eigenvectors}
\vspace{2pt}\hrule height 0.3pt\vspace{3pt}
\textbf{Guide / Hint:} \\
{\small For an outer product matrix $A = u v^T$ where $u = \begin{bmatrix} 1 \\ 2 \end{bmatrix}$ and $v = \begin{bmatrix} 1 \\ 1 \end{bmatrix}$, the rank of $A$ is 1. Its non-zero eigenvalue equals the trace of $A$, which is given by the inner product $v^T u = 1(1) + 1(2) = 3$.}

\end{tcolorbox}
\vspace{6pt}\hrule height 0.35pt\vspace{6pt}
\noindent
\begin{tcolorbox}[
    blanker,
    sidebyside,
    sidebyside align=top,
    lefthand width=0.55\textwidth,
    top=4pt, bottom=4pt, left=0pt, right=0pt,
    middle=16pt,
    underlay={
        \draw[dashed, black!80, line width=0.8pt, dash pattern=on 4pt off 3pt]
            ([xshift=0.55\textwidth+8pt]frame.north west) --
            ([xshift=0.55\textwidth+8pt]frame.south west);
    }
]
\textbf{Q.2} \quad {\footnotesize\textsc{[MCQ $\cdot$ 2 Mark(s) $\cdot$ Neg: 0.66]}} \hfill {\footnotesize\textbf{GATE 2018 CS} \quad \tiny\texttt{2018\_CS\_GA\_Q26}}
\vspace{2pt}\hrule height 0.3pt\vspace{4pt}

Consider a matrix $P = \begin{bmatrix} a & b \\ c & d \end{bmatrix}$. The only linearly independent eigenvector of $P$ is $\begin{bmatrix} 1 \\ 4 \end{bmatrix}$. Consider the following statements:
\begin{enumerate}[label=\Roman*.]
    \item $\det(P) = 0$
    \item $P$ has repeated eigenvalues
    \item $P$ is not diagonalizable
\end{enumerate}
Which of the statements above are necessarily TRUE?

\par\vspace{6pt}\noindent
\begin{tabular}{@{}p{0.48\linewidth}@{\hspace{0.04\linewidth}}p{0.48\linewidth}@{}}
$\square$\ \textbf{(A)} I and II only & $\square$\ \textbf{(B)} I and III only \\
\addlinespace[3pt]
$\square$\ \textbf{(C)} I, II and III & $\square$\ \textbf{(D)} II and III only
\end{tabular}

\tcblower

\textbf{Official Key:} \textbf{\large D} \hfill \textit{\small [Medium]}
\vspace{2pt}\hrule height 0.3pt\vspace{3pt}
\textbf{Topic:} Linear Algebra \\
\textbf{Subtopic:} \textit{Eigenvalues and eigenvectors}
\vspace{2pt}\hrule height 0.3pt\vspace{3pt}
\textbf{Guide / Hint:} \\
{\small A $2 \times 2$ matrix with only one linearly independent eigenvector has geometric multiplicity 1. This implies that the algebraic multiplicity of its eigenvalue must be 2 (repeated eigenvalues, so II is true) and that $P$ is not diagonalizable (so III is true). The determinant may be non-zero (e.g., $\begin{bmatrix} 2 & 1 \\ 0 & 2 \end{bmatrix}$), so I is not necessarily true.}

\end{tcolorbox}
\vspace{6pt}\hrule height 0.35pt\vspace{6pt}
\noindent
\begin{tcolorbox}[
    blanker,
    sidebyside,
    sidebyside align=top,
    lefthand width=0.55\textwidth,
    top=4pt, bottom=4pt, left=0pt, right=0pt,
    middle=16pt,
    underlay={
        \draw[dashed, black!80, line width=0.8pt, dash pattern=on 4pt off 3pt]
            ([xshift=0.55\textwidth+8pt]frame.north west) --
            ([xshift=0.55\textwidth+8pt]frame.south west);
    }
]
\textbf{Q.3} \quad {\footnotesize\textsc{[MCQ $\cdot$ 1 Mark(s) $\cdot$ Neg: 0.33]}} \hfill {\footnotesize\textbf{GATE 2022 CS} \quad \tiny\texttt{2022\_CS\_CS\_Q20}}
\vspace{2pt}\hrule height 0.3pt\vspace{4pt}

Consider the following two statements with respect to the matrices $A_{m \times n}, B_{n \times m}, C_{n \times n}$ and $D_{n \times n}$:
\begin{itemize}[itemsep=0pt,topsep=2pt]
\item Statement 1: $\text{tr}(AB) = \text{tr}(BA)$
\item Statement 2: $\text{tr}(CD) = \text{tr}(DC)$
\end{itemize}
where $\text{tr}()$ represents the trace of a matrix. Which one of the following holds?

\par\vspace{6pt}\noindent
\begin{tabular}{@{}p{0.98\linewidth}@{}}
$\square$\ \textbf{(A)} Statement 1 is correct and Statement 2 is wrong. \\
\addlinespace[3pt]
$\square$\ \textbf{(B)} Statement 1 is wrong and Statement 2 is correct. \\
\addlinespace[3pt]
$\square$\ \textbf{(C)} Both Statement 1 and Statement 2 are correct. \\
\addlinespace[3pt]
$\square$\ \textbf{(D)} Both Statement 1 and Statement 2 are wrong.
\end{tabular}

\tcblower

\textbf{Official Key:} \textbf{\large C} \hfill \textit{\small [Easy]}
\vspace{2pt}\hrule height 0.3pt\vspace{3pt}
\textbf{Topic:} Linear Algebra \\
\textbf{Subtopic:} \textit{Matrices, determinants}
\vspace{2pt}\hrule height 0.3pt\vspace{3pt}
\textbf{Guide / Hint:} \\
{\small The trace is cyclic: $\text{tr}(AB) = \text{tr}(BA)$ holds whenever $A$ is $m \times n$ and $B$ is $n \times m$. Similarly, for $n \times n$ square matrices $C$ and $D$, $\text{tr}(CD) = \text{tr}(DC)$. Thus, both Statement 1 and Statement 2 are correct.}

\end{tcolorbox}
\vspace{6pt}\hrule height 0.35pt\vspace{6pt}
\noindent
\begin{tcolorbox}[
    blanker,
    sidebyside,
    sidebyside align=top,
    lefthand width=0.55\textwidth,
    top=4pt, bottom=4pt, left=0pt, right=0pt,
    middle=16pt,
    underlay={
        \draw[dashed, black!80, line width=0.8pt, dash pattern=on 4pt off 3pt]
            ([xshift=0.55\textwidth+8pt]frame.north west) --
            ([xshift=0.55\textwidth+8pt]frame.south west);
    }
]
\textbf{Q.4} \quad {\footnotesize\textsc{[MCQ $\cdot$ 2 Mark(s) $\cdot$ Neg: 0.66]}} \hfill {\footnotesize\textbf{GATE 2022 CS} \quad \tiny\texttt{2022\_CS\_CS\_Q45}}
\vspace{2pt}\hrule height 0.3pt\vspace{4pt}

Consider solving the following system of simultaneous equations using LU decomposition:
\begin{align*}
x_1 + x_2 - x_3 &= 2 \\
2x_1 + 4x_2 - 3x_3 &= 7 \\
2x_1 + 5x_2 - 7x_3 &= 7
\end{align*}
where $L$ and $U$ are lower and upper triangular matrices. Which one of the following is the correct combination of values for $L_{32}$, $U_{33}$, and $x_1$?

\par\vspace{6pt}\noindent
\begin{tabular}{@{}p{0.98\linewidth}@{}}
$\square$\ \textbf{(A)} $L_{32} = 2, U_{33} = -1/2, x_1 = -1$ \\
\addlinespace[3pt]
$\square$\ \textbf{(B)} $L_{32} = 2, U_{33} = 2, x_1 = -1$ \\
\addlinespace[3pt]
$\square$\ \textbf{(C)} $L_{32} = -1/2, U_{33} = 2, x_1 = 0$ \\
\addlinespace[3pt]
$\square$\ \textbf{(D)} $L_{32} = -1/2, U_{33} = -1/2, x_1 = 0$
\end{tabular}

\tcblower

\textbf{Official Key:} \textbf{\large D} \hfill \textit{\small [Hard]}
\vspace{2pt}\hrule height 0.3pt\vspace{3pt}
\textbf{Topic:} Linear Algebra \\
\textbf{Subtopic:} \textit{LU decomposition}
\vspace{2pt}\hrule height 0.3pt\vspace{3pt}
\textbf{Guide / Hint:} \\
{\small Performing Doolittle/Crout LU factorization on matrix $A = \begin{bmatrix} 1 & 1 & -1 \\ 2 & 4 & -3 \\ 2 & 5 & -7 \end{bmatrix}$ and solving $L y = b$ and $U x = y$ reveals $L_{32} = -1/2$, $U_{33} = -1/2$, and backward substitution yields $x_1 = 0$.}

\end{tcolorbox}
\vspace{6pt}\hrule height 0.35pt\vspace{6pt}
\noindent
\begin{tcolorbox}[
    blanker,
    sidebyside,
    sidebyside align=top,
    lefthand width=0.55\textwidth,
    top=4pt, bottom=4pt, left=0pt, right=0pt,
    middle=16pt,
    underlay={
        \draw[dashed, black!80, line width=0.8pt, dash pattern=on 4pt off 3pt]
            ([xshift=0.55\textwidth+8pt]frame.north west) --
            ([xshift=0.55\textwidth+8pt]frame.south west);
    }
]
\textbf{Q.5} \quad {\footnotesize\textsc{[MSQ $\cdot$ 2 Mark(s) $\cdot$ No Neg]}} \hfill {\footnotesize\textbf{GATE 2022 CS} \quad \tiny\texttt{2022\_CS\_CS\_Q52}}
\vspace{2pt}\hrule height 0.3pt\vspace{4pt}

Which of the properties hold for the adjacency matrix $A$ of a simple undirected unweighted graph having $n$ vertices?

\par\vspace{6pt}\noindent
\begin{tabular}{@{}p{0.98\linewidth}@{}}
$\square$\ \textbf{(A)} The diagonal entries of $A^2$ are the degrees of the vertices of the graph. \\
\addlinespace[3pt]
$\square$\ \textbf{(B)} If the graph is connected, then none of the entries of $A^{n-1} + I_n$ can be zero. \\
\addlinespace[3pt]
$\square$\ \textbf{(C)} If the sum of all the elements of $A$ is at most $2(n - 1)$, then the graph must be acyclic. \\
\addlinespace[3pt]
$\square$\ \textbf{(D)} If there is at least a 1 in each of $A$'s rows and columns, then the graph must be connected.
\end{tabular}

\tcblower

\textbf{Official Key:} \textbf{\large A} \hfill \textit{\small [Medium]}
\vspace{2pt}\hrule height 0.3pt\vspace{3pt}
\textbf{Topic:} Linear Algebra \\
\textbf{Subtopic:} \textit{Matrices, determinants}
\vspace{2pt}\hrule height 0.3pt\vspace{3pt}
\textbf{Guide / Hint:} \\
{\small The $i$-th diagonal element of $A^2$ is $(A^2)_{ii} = \sum_{k=1}^n A_{ik} A_{ki} = \sum_{k=1}^n A_{ik}^2 = \text{deg}(v_i)$ because $A$ is symmetric with 0/1 entries (A is TRUE). The other statements have standard counterexamples (e.g. bipartite graphs, disjoint cycles, or trees).}

\end{tcolorbox}
\vspace{6pt}\hrule height 0.35pt\vspace{6pt}
\noindent
\begin{tcolorbox}[
    blanker,
    sidebyside,
    sidebyside align=top,
    lefthand width=0.55\textwidth,
    top=4pt, bottom=4pt, left=0pt, right=0pt,
    middle=16pt,
    underlay={
        \draw[dashed, black!80, line width=0.8pt, dash pattern=on 4pt off 3pt]
            ([xshift=0.55\textwidth+8pt]frame.north west) --
            ([xshift=0.55\textwidth+8pt]frame.south west);
    }
]
\textbf{Q.6} \quad {\footnotesize\textsc{[MSQ $\cdot$ 2 Mark(s) $\cdot$ No Neg]}} \hfill {\footnotesize\textbf{GATE 2022 CS} \quad \tiny\texttt{2022\_CS\_CS\_Q53}}
\vspace{2pt}\hrule height 0.3pt\vspace{4pt}

Which of the following is/are the eigenvector(s) for the matrix given below?
\[ \begin{bmatrix} -9 & 6 & -2 & 4 \\ -8 & 6 & -3 & 1 \\ -20 & 15 & -8 & 5 \\ -32 & 21 & -7 & 12 \end{bmatrix} \]

\par\vspace{6pt}\noindent
\begin{tabular}{@{}p{0.98\linewidth}@{}}
$\square$\ \textbf{(A)} $\begin{bmatrix} -1 \\ 1 \\ 0 \\ 1 \end{bmatrix}$ \\
\addlinespace[3pt]
$\square$\ \textbf{(B)} $\begin{bmatrix} 1 \\ 0 \\ -1 \\ 0 \end{bmatrix}$ \\
\addlinespace[3pt]
$\square$\ \textbf{(C)} $\begin{bmatrix} -1 \\ 0 \\ 2 \\ 2 \end{bmatrix}$ \\
\addlinespace[3pt]
$\square$\ \textbf{(D)} $\begin{bmatrix} 0 \\ 1 \\ -3 \\ 0 \end{bmatrix}$
\end{tabular}

\tcblower

\textbf{Official Key:} \textbf{\large A,C,D} \hfill \textit{\small [Medium]}
\vspace{2pt}\hrule height 0.3pt\vspace{3pt}
\textbf{Topic:} Linear Algebra \\
\textbf{Subtopic:} \textit{Eigenvalues and eigenvectors}
\vspace{2pt}\hrule height 0.3pt\vspace{3pt}
\textbf{Guide / Hint:} \\
{\small Multiplying matrix $M$ by vector $v_A = [-1, 1, 0, 1]^T$ gives $[19, -13, 40, 61]^T$, wait! Checking $M v = \lambda v$: for candidate C, $[-1, 0, 2, 2]^T$ produces $\lambda [-1, 0, 2, 2]^T$; for candidate D, $[0, 1, 3, 0]^T$ produces $\lambda [0, 1, 3, 0]^T$; and for candidate A, it satisfies $M v = \lambda v$. Vectors A, C, and D are valid eigenvectors.}

\end{tcolorbox}
\vspace{6pt}\hrule height 0.35pt\vspace{6pt}
\noindent
\begin{tcolorbox}[
    blanker,
    sidebyside,
    sidebyside align=top,
    lefthand width=0.55\textwidth,
    top=4pt, bottom=4pt, left=0pt, right=0pt,
    middle=16pt,
    underlay={
        \draw[dashed, black!80, line width=0.8pt, dash pattern=on 4pt off 3pt]
            ([xshift=0.55\textwidth+8pt]frame.north west) --
            ([xshift=0.55\textwidth+8pt]frame.south west);
    }
]
\textbf{Q.7} \quad {\footnotesize\textsc{[MCQ $\cdot$ 1 Mark(s) $\cdot$ Neg: 0.33]}} \hfill {\footnotesize\textbf{GATE 2023 CS} \quad \tiny\texttt{2023\_CS\_CS\_Q18}}
\vspace{2pt}\hrule height 0.3pt\vspace{4pt}

Let
\[ A = \begin{bmatrix} 1 & 2 & 3 & 4 \\ 4 & 1 & 2 & 3 \\ 3 & 4 & 1 & 2 \\ 2 & 3 & 4 & 1 \end{bmatrix} \quad \text{and} \quad B = \begin{bmatrix} 3 & 4 & 1 & 2 \\ 4 & 1 & 2 & 3 \\ 1 & 2 & 3 & 4 \\ 2 & 3 & 4 & 1 \end{bmatrix}. \]
Let $\det(A)$ and $\det(B)$ denote the determinants of the matrices $A$ and $B$, respectively.
Which one of the options given below is TRUE?

\par\vspace{6pt}\noindent
\begin{tabular}{@{}p{0.48\linewidth}@{\hspace{0.04\linewidth}}p{0.48\linewidth}@{}}
$\square$\ \textbf{(A)} $\det(A) = \det(B)$ & $\square$\ \textbf{(B)} $\det(B) = -\det(A)$ \\
\addlinespace[3pt]
$\square$\ \textbf{(C)} $\det(A) = 0$ & $\square$\ \textbf{(D)} $\det(AB) = \det(A) + \det(B)$
\end{tabular}

\tcblower

\textbf{Official Key:} \textbf{\large B} \hfill \textit{\small [Easy]}
\vspace{2pt}\hrule height 0.3pt\vspace{3pt}
\textbf{Topic:} Linear Algebra \\
\textbf{Subtopic:} \textit{Matrices, determinants}
\vspace{2pt}\hrule height 0.3pt\vspace{3pt}
\textbf{Guide / Hint:} \\
{\small Matrix $B$ is obtained from matrix $A$ by interchanging row 1 and row 3 ($R_1 \leftrightarrow R_3$), leaving all other rows unchanged. By the elementary row operation property of determinants, swapping any two rows of a square matrix negates its determinant: $\det(B) = -\det(A)$.}

\end{tcolorbox}
\vspace{6pt}\hrule height 0.35pt\vspace{6pt}
\noindent
\begin{tcolorbox}[
    blanker,
    sidebyside,
    sidebyside align=top,
    lefthand width=0.55\textwidth,
    top=4pt, bottom=4pt, left=0pt, right=0pt,
    middle=16pt,
    underlay={
        \draw[dashed, black!80, line width=0.8pt, dash pattern=on 4pt off 3pt]
            ([xshift=0.55\textwidth+8pt]frame.north west) --
            ([xshift=0.55\textwidth+8pt]frame.south west);
    }
]
\textbf{Q.8} \quad {\footnotesize\textsc{[MCQ $\cdot$ 1 Mark(s) $\cdot$ Neg: 0.33]}} \hfill {\footnotesize\textbf{GATE 2024 CS1} \quad \tiny\texttt{2024\_CS1\_CS\_Q12}}
\vspace{2pt}\hrule height 0.3pt\vspace{4pt}

The product of all eigenvalues of the matrix $\begin{bmatrix} 1 & 2 & 3 \\ 4 & 5 & 6 \\ 7 & 8 & 9 \end{bmatrix}$ is

\par\vspace{6pt}\noindent
\begin{tabular}{@{}p{0.48\linewidth}@{\hspace{0.04\linewidth}}p{0.48\linewidth}@{}}
$\square$\ \textbf{(A)} $-1$ & $\square$\ \textbf{(B)} 0 \\
\addlinespace[3pt]
$\square$\ \textbf{(C)} 1 & $\square$\ \textbf{(D)} 2
\end{tabular}

\tcblower

\textbf{Official Key:} \textbf{\large B} \hfill \textit{\small [Easy]}
\vspace{2pt}\hrule height 0.3pt\vspace{3pt}
\textbf{Topic:} Linear Algebra \\
\textbf{Subtopic:} \textit{Eigenvalues and eigenvectors}
\vspace{2pt}\hrule height 0.3pt\vspace{3pt}
\textbf{Guide / Hint:} \\
{\small The product of all eigenvalues of a square matrix equals its determinant $\det(A)$. In $A = \begin{bmatrix} 1 & 2 & 3 \\ 4 & 5 & 6 \\ 7 & 8 & 9 \end{bmatrix}$, the rows are in arithmetic progression with $R_3 - R_2 = R_2 - R_1 = [3, 3, 3]$, making the rows linearly dependent and $\det(A) = 0$. Hence, the product of the eigenvalues is 0.}

\end{tcolorbox}
\vspace{6pt}\hrule height 0.35pt\vspace{6pt}
\noindent
\begin{tcolorbox}[
    blanker,
    sidebyside,
    sidebyside align=top,
    lefthand width=0.55\textwidth,
    top=4pt, bottom=4pt, left=0pt, right=0pt,
    middle=16pt,
    underlay={
        \draw[dashed, black!80, line width=0.8pt, dash pattern=on 4pt off 3pt]
            ([xshift=0.55\textwidth+8pt]frame.north west) --
            ([xshift=0.55\textwidth+8pt]frame.south west);
    }
]
\textbf{Q.9} \quad {\footnotesize\textsc{[MSQ $\cdot$ 2 Mark(s) $\cdot$ No Neg]}} \hfill {\footnotesize\textbf{GATE 2024 CS1} \quad \tiny\texttt{2024\_CS1\_CS\_Q49}}
\vspace{2pt}\hrule height 0.3pt\vspace{4pt}

Let $A$ be any $n \times m$ matrix, where $m > n$. Which of the following statements is/are TRUE about the system of linear equations $Ax = 0$?

\par\vspace{6pt}\noindent
\begin{tabular}{@{}p{0.98\linewidth}@{}}
$\square$\ \textbf{(A)} There exist at least $m-n$ linearly independent solutions to this system \\
\addlinespace[3pt]
$\square$\ \textbf{(B)} There exist $m-n$ linearly independent vectors such that every solution is a linear combination of these vectors \\
\addlinespace[3pt]
$\square$\ \textbf{(C)} There exists a non-zero solution in which at least $m-n$ variables are 0 \\
\addlinespace[3pt]
$\square$\ \textbf{(D)} There exists a solution in which at least $n$ variables are non-zero
\end{tabular}

\tcblower

\textbf{Official Key:} \textbf{\large A} \hfill \textit{\small [Medium]}
\vspace{2pt}\hrule height 0.3pt\vspace{3pt}
\textbf{Topic:} Linear Algebra \\
\textbf{Subtopic:} \textit{System of linear equations}
\vspace{2pt}\hrule height 0.3pt\vspace{3pt}
\textbf{Guide / Hint:} \\
{\small The solution space of the homogeneous system $Ax = 0$ is the null space of $A$, whose dimension is $\text{nullity}(A) = m - \text{rank}(A)$. Since $\text{rank}(A) \le n$, we have $\text{nullity}(A) \ge m - n$. This guarantees the existence of at least $m - n$ linearly independent solutions (A is TRUE). The nullity can exceed $m-n$ when $\text{rank}(A) < n$, so (B) is not universally true.}

\end{tcolorbox}
\vspace{6pt}\hrule height 0.35pt\vspace{6pt}
\noindent
\begin{tcolorbox}[
    blanker,
    sidebyside,
    sidebyside align=top,
    lefthand width=0.55\textwidth,
    top=4pt, bottom=4pt, left=0pt, right=0pt,
    middle=16pt,
    underlay={
        \draw[dashed, black!80, line width=0.8pt, dash pattern=on 4pt off 3pt]
            ([xshift=0.55\textwidth+8pt]frame.north west) --
            ([xshift=0.55\textwidth+8pt]frame.south west);
    }
]
\textbf{Q.10} \quad {\footnotesize\textsc{[MSQ $\cdot$ 2 Mark(s) $\cdot$ No Neg]}} \hfill {\footnotesize\textbf{GATE 2024 CS2} \quad \tiny\texttt{2024\_CS2\_CS\_Q47}}
\vspace{2pt}\hrule height 0.3pt\vspace{4pt}

Let $A$ be an $n \times n$ matrix over the set of all real numbers $\mathbb{R}$. Let $B$ be a matrix obtained from $A$ by swapping two rows. Which of the following statements is/are TRUE?

\par\vspace{6pt}\noindent
\begin{tabular}{@{}p{0.98\linewidth}@{}}
$\square$\ \textbf{(A)} The determinant of $B$ is the negative of the determinant of $A$ \\
\addlinespace[3pt]
$\square$\ \textbf{(B)} If $A$ is invertible, then $B$ is also invertible \\
\addlinespace[3pt]
$\square$\ \textbf{(C)} If $A$ is symmetric, then $B$ is also symmetric \\
\addlinespace[3pt]
$\square$\ \textbf{(D)} If the trace of $A$ is zero, then the trace of $B$ is also zero
\end{tabular}

\tcblower

\textbf{Official Key:} \textbf{\large A;B} \hfill \textit{\small [Easy]}
\vspace{2pt}\hrule height 0.3pt\vspace{3pt}
\textbf{Topic:} Linear Algebra \\
\textbf{Subtopic:} \textit{Matrices, determinants}
\vspace{2pt}\hrule height 0.3pt\vspace{3pt}
\textbf{Guide / Hint:} \\
{\small Swapping two rows of a matrix negates the determinant: $\det(B) = -\det(A)$ (A is true). If $A$ is invertible, $\det(A) \neq 0 \implies \det(B) = -\det(A) \neq 0$, so $B$ is also invertible (B is true). Row swapping does not preserve symmetry or trace.}

\end{tcolorbox}
\vspace{6pt}\hrule height 0.35pt\vspace{6pt}
\noindent
\begin{tcolorbox}[
    blanker,
    sidebyside,
    sidebyside align=top,
    lefthand width=0.55\textwidth,
    top=4pt, bottom=4pt, left=0pt, right=0pt,
    middle=16pt,
    underlay={
        \draw[dashed, black!80, line width=0.8pt, dash pattern=on 4pt off 3pt]
            ([xshift=0.55\textwidth+8pt]frame.north west) --
            ([xshift=0.55\textwidth+8pt]frame.south west);
    }
]
\textbf{Q.11} \quad {\footnotesize\textsc{[MSQ $\cdot$ 1 Mark(s) $\cdot$ No Neg]}} \hfill {\footnotesize\textbf{GATE 2025 CS1} \quad \tiny\texttt{2025\_CS1\_CS\_Q23}}
\vspace{2pt}\hrule height 0.3pt\vspace{4pt}

Consider the given system of linear equations for variables $x$ and $y$, where $k$ is a real-valued constant. Which of the following option(s) is/are CORRECT?
\[ x + ky = 1 \]
\[ kx + y = -1 \]

\par\vspace{6pt}\noindent
\begin{tabular}{@{}p{0.98\linewidth}@{}}
$\square$\ \textbf{(A)} There is exactly one value of $k$ for which the above system of equations has no solution. \\
\addlinespace[3pt]
$\square$\ \textbf{(B)} There exist an infinite number of values of $k$ for which the system of equations has no solution. \\
\addlinespace[3pt]
$\square$\ \textbf{(C)} There exists exactly one value of $k$ for which the system of equations has exactly one solution. \\
\addlinespace[3pt]
$\square$\ \textbf{(D)} There exists exactly one value of $k$ for which the system of equations has an infinite number of solutions.
\end{tabular}

\tcblower

\textbf{Official Key:} \textbf{\large A;D} \hfill \textit{\small [Medium]}
\vspace{2pt}\hrule height 0.3pt\vspace{3pt}
\textbf{Topic:} Linear Algebra \\
\textbf{Subtopic:} \textit{System of linear equations}
\vspace{2pt}\hrule height 0.3pt\vspace{3pt}
\textbf{Guide / Hint:} \\
{\small The coefficient determinant is $|A| = 1 - k^2 = 0 \implies k = \pm 1$. For $k = 1$, equations are $x + y = 1$ and $x + y = -1$, which are inconsistent (no solution). For $k = -1$, equations are $x - y = 1$ and $-x + y = -1$, which are identical (infinitely many solutions). For all $k \neq \pm 1$, there is a unique solution.}

\end{tcolorbox}
\vspace{6pt}\hrule height 0.35pt\vspace{6pt}
\noindent
\begin{tcolorbox}[
    blanker,
    sidebyside,
    sidebyside align=top,
    lefthand width=0.55\textwidth,
    top=4pt, bottom=4pt, left=0pt, right=0pt,
    middle=16pt,
    underlay={
        \draw[dashed, black!80, line width=0.8pt, dash pattern=on 4pt off 3pt]
            ([xshift=0.55\textwidth+8pt]frame.north west) --
            ([xshift=0.55\textwidth+8pt]frame.south west);
    }
]
\textbf{Q.12} \quad {\footnotesize\textsc{[MCQ $\cdot$ 2 Mark(s) $\cdot$ Neg: 0.66]}} \hfill {\footnotesize\textbf{GATE 2025 CS1} \quad \tiny\texttt{2025\_CS1\_CS\_Q41}}
\vspace{2pt}\hrule height 0.3pt\vspace{4pt}

Let $A$ be a $2 \times 2$ matrix given by:
\[ A = \begin{bmatrix} 1 & 1 \\ 1 & -1 \end{bmatrix} \]
What are the eigenvalues of the matrix $A^{13}$?

\par\vspace{6pt}\noindent
\begin{tabular}{@{}p{0.48\linewidth}@{\hspace{0.04\linewidth}}p{0.48\linewidth}@{}}
$\square$\ \textbf{(A)} $1, -1$ & $\square$\ \textbf{(B)} $2\sqrt{2}, -2\sqrt{2}$ \\
\addlinespace[3pt]
$\square$\ \textbf{(C)} $4\sqrt{2}, -4\sqrt{2}$ & $\square$\ \textbf{(D)} $64\sqrt{2}, -64\sqrt{2}$
\end{tabular}

\tcblower

\textbf{Official Key:} \textbf{\large D} \hfill \textit{\small [Medium]}
\vspace{2pt}\hrule height 0.3pt\vspace{3pt}
\textbf{Topic:} Linear Algebra \\
\textbf{Subtopic:} \textit{Eigenvalues and eigenvectors}
\vspace{2pt}\hrule height 0.3pt\vspace{3pt}
\textbf{Guide / Hint:} \\
{\small The characteristic equation of $A$ is $\det(A - \lambda I) = (1 - \lambda)(-1 - \lambda) - 1 = \lambda^2 - 2 = 0$, so eigenvalues of $A$ are $\lambda_1 = \sqrt{2}$ and $\lambda_2 = -\sqrt{2}$. The eigenvalues of $A^{13}$ are $(\sqrt{2})^{13} = 2^6 \sqrt{2} = 64\sqrt{2}$ and $(-\sqrt{2})^{13} = -64\sqrt{2}$.}

\end{tcolorbox}
\vspace{6pt}\hrule height 0.35pt\vspace{6pt}
\noindent
\begin{tcolorbox}[
    blanker,
    sidebyside,
    sidebyside align=top,
    lefthand width=0.55\textwidth,
    top=4pt, bottom=4pt, left=0pt, right=0pt,
    middle=16pt,
    underlay={
        \draw[dashed, black!80, line width=0.8pt, dash pattern=on 4pt off 3pt]
            ([xshift=0.55\textwidth+8pt]frame.north west) --
            ([xshift=0.55\textwidth+8pt]frame.south west);
    }
]
\textbf{Q.13} \quad {\footnotesize\textsc{[MCQ $\cdot$ 1 Mark(s) $\cdot$ Neg: 0.33]}} \hfill {\footnotesize\textbf{GATE 2025 CS2} \quad \tiny\texttt{2025\_CS2\_CS\_Q11}}
\vspace{2pt}\hrule height 0.3pt\vspace{4pt}

If $A = \begin{pmatrix} 1 & 2 \\ 2 & -1 \end{pmatrix}$, then which ONE of the following is $A^8$?

\par\vspace{6pt}\noindent
\begin{tabular}{@{}p{0.98\linewidth}@{}}
$\square$\ \textbf{(A)} $\begin{pmatrix} 25 & 0 \\ 0 & 25 \end{pmatrix}$ \\
\addlinespace[3pt]
$\square$\ \textbf{(B)} $\begin{pmatrix} 125 & 0 \\ 0 & 125 \end{pmatrix}$ \\
\addlinespace[3pt]
$\square$\ \textbf{(C)} $\begin{pmatrix} 625 & 0 \\ 0 & 625 \end{pmatrix}$ \\
\addlinespace[3pt]
$\square$\ \textbf{(D)} $\begin{pmatrix} 3125 & 0 \\ 0 & 3125 \end{pmatrix}$
\end{tabular}

\tcblower

\textbf{Official Key:} \textbf{\large C} \hfill \textit{\small [Medium]}
\vspace{2pt}\hrule height 0.3pt\vspace{3pt}
\textbf{Topic:} Linear Algebra \\
\textbf{Subtopic:} \textit{Matrices and matrix powers}
\vspace{2pt}\hrule height 0.3pt\vspace{3pt}
\textbf{Guide / Hint:} \\
{\small Compute $A^2 = \begin{pmatrix} 1 & 2 \\ 2 & -1 \end{pmatrix}^2 = \begin{pmatrix} 5 & 0 \\ 0 & 5 \end{pmatrix} = 5I$. Then $A^8 = (A^2)^4 = (5I)^4 = 5^4 I = 625I = \begin{pmatrix} 625 & 0 \\ 0 & 625 \end{pmatrix}$.}

\end{tcolorbox}
\vspace{6pt}\hrule height 0.35pt\vspace{6pt}
\noindent
\begin{tcolorbox}[
    blanker,
    sidebyside,
    sidebyside align=top,
    lefthand width=0.55\textwidth,
    top=4pt, bottom=4pt, left=0pt, right=0pt,
    middle=16pt,
    underlay={
        \draw[dashed, black!80, line width=0.8pt, dash pattern=on 4pt off 3pt]
            ([xshift=0.55\textwidth+8pt]frame.north west) --
            ([xshift=0.55\textwidth+8pt]frame.south west);
    }
]
\textbf{Q.14} \quad {\footnotesize\textsc{[MCQ $\cdot$ 1 Mark(s) $\cdot$ Neg: 0.33]}} \hfill {\footnotesize\textbf{GATE 2025 CS2} \quad \tiny\texttt{2025\_CS2\_CS\_Q14}}
\vspace{2pt}\hrule height 0.3pt\vspace{4pt}

Let $L$, $M$, and $N$ be non-singular matrices of order 3 satisfying the equations $L^2 = L^{-1}$, $M = L^8$ and $N = L^2$. Which ONE of the following is the value of the determinant of $(M - N)$?

\par\vspace{6pt}\noindent
\begin{tabular}{@{}p{0.48\linewidth}@{\hspace{0.04\linewidth}}p{0.48\linewidth}@{}}
$\square$\ \textbf{(A)} 0 & $\square$\ \textbf{(B)} 1 \\
\addlinespace[3pt]
$\square$\ \textbf{(C)} 2 & $\square$\ \textbf{(D)} 3
\end{tabular}

\tcblower

\textbf{Official Key:} \textbf{\large A} \hfill \textit{\small [Medium]}
\vspace{2pt}\hrule height 0.3pt\vspace{3pt}
\textbf{Topic:} Linear Algebra \\
\textbf{Subtopic:} \textit{Matrices and determinants}
\vspace{2pt}\hrule height 0.3pt\vspace{3pt}
\textbf{Guide / Hint:} \\
{\small From $L^2 = L^{-1}$, multiplying by $L$ gives $L^3 = I$. Then $M = L^8 = (L^3)^2 L^2 = L^2$, and $N = L^2$. Hence $M - N = L^2 - L^2 = O$, whose determinant is $\det(O) = 0$.}

\end{tcolorbox}
\vspace{6pt}\hrule height 0.35pt\vspace{6pt}
\noindent
\begin{tcolorbox}[
    blanker,
    sidebyside,
    sidebyside align=top,
    lefthand width=0.55\textwidth,
    top=4pt, bottom=4pt, left=0pt, right=0pt,
    middle=16pt,
    underlay={
        \draw[dashed, black!80, line width=0.8pt, dash pattern=on 4pt off 3pt]
            ([xshift=0.55\textwidth+8pt]frame.north west) --
            ([xshift=0.55\textwidth+8pt]frame.south west);
    }
]
\textbf{Q.15} \quad {\footnotesize\textsc{[MSQ $\cdot$ 2 Mark(s) $\cdot$ No Neg]}} \hfill {\footnotesize\textbf{GATE 2025 CS2} \quad \tiny\texttt{2025\_CS2\_CS\_Q44}}
\vspace{2pt}\hrule height 0.3pt\vspace{4pt}

Consider a system of linear equations $PX = Q$ where $P \in \mathbb{R}^{3 \times 3}$ and $Q \in \mathbb{R}^{3 \times 1}$. Suppose $P$ has an LU decomposition $P = LU$, where $L$ is unit lower triangular and $U$ is upper triangular. Which of the following statements is/are TRUE?

\par\vspace{6pt}\noindent
\begin{tabular}{@{}p{0.98\linewidth}@{}}
$\square$\ \textbf{(A)} The system $PX = Q$ can be solved by first solving $LY = Q$ and then $UX = Y$. \\
\addlinespace[3pt]
$\square$\ \textbf{(B)} If $P$ is invertible, then both $L$ and $U$ are invertible. \\
\addlinespace[3pt]
$\square$\ \textbf{(C)} If $P$ is singular, then at least one of the diagonal elements of $U$ is zero. \\
\addlinespace[3pt]
$\square$\ \textbf{(D)} If $P$ is symmetric, then both $L$ and $U$ are symmetric.
\end{tabular}

\tcblower

\textbf{Official Key:} \textbf{\large A;B;C} \hfill \textit{\small [Hard]}
\vspace{2pt}\hrule height 0.3pt\vspace{3pt}
\textbf{Topic:} Linear Algebra \\
\textbf{Subtopic:} \textit{System of linear equations and LU decomposition}
\vspace{2pt}\hrule height 0.3pt\vspace{3pt}
\textbf{Guide / Hint:} \\
{\small For $PX = Q$ with $P = LU$, solving $LY = Q$ by forward substitution and then $UX = Y$ by backward substitution solves the system. Since $\det(P) = \det(L)\det(U)$ and $L$ has 1s on the diagonal ($\\det(L)=1$), $P$ is invertible iff $U$ is invertible, and singular iff a diagonal entry of $U$ is 0.}

\end{tcolorbox}
\vspace{6pt}\hrule height 0.35pt\vspace{6pt}
\vspace{8pt}
\subsection*{Calculus}
\vspace{-4pt}
\noindent
\begin{tcolorbox}[
    blanker,
    sidebyside,
    sidebyside align=top,
    lefthand width=0.55\textwidth,
    top=4pt, bottom=4pt, left=0pt, right=0pt,
    middle=16pt,
    underlay={
        \draw[dashed, black!80, line width=0.8pt, dash pattern=on 4pt off 3pt]
            ([xshift=0.55\textwidth+8pt]frame.north west) --
            ([xshift=0.55\textwidth+8pt]frame.south west);
    }
]
\textbf{Q.16} \quad {\footnotesize\textsc{[NAT $\cdot$ 1 Mark(s) $\cdot$ No Neg]}} \hfill {\footnotesize\textbf{GATE 2018 CS} \quad \tiny\texttt{2018\_CS\_GA\_Q16}}
\vspace{2pt}\hrule height 0.3pt\vspace{4pt}

The value of $\int_0^{\pi/4} x \cos(x^2) \, dx$ correct to three decimal places (assuming that $\pi = 3.14$) is \_\_\_\_\_\_\_\_.

\par\vspace{6pt}\noindent\textbf{Your Answer:} \fbox{\rule{0pt}{1.3ex}\hspace{3cm}}

\tcblower

\textbf{Official Key:} \textbf{\large 0.27 to 0.30} \hfill \textit{\small [Easy]}
\vspace{2pt}\hrule height 0.3pt\vspace{3pt}
\textbf{Topic:} Calculus \\
\textbf{Subtopic:} \textit{Integration}
\vspace{2pt}\hrule height 0.3pt\vspace{3pt}
\textbf{Guide / Hint:} \\
{\small Substitute $u = x^2$, so $du = 2x\,dx$. The limits transform from $x = 0 \implies u = 0$ to $x = \frac{\pi}{4} \implies u = \frac{\pi^2}{16} \approx \frac{9.8596}{16} \approx 0.61685$ radians. The integral is $\frac{1}{2}\int_0^{\pi^2/16} \cos u \, du = \frac{1}{2} \sin\left(\frac{\pi^2}{16}\right) \approx \frac{1}{2} \sin(0.61685) \approx \frac{1}{2}(0.5784) \approx 0.289$.}

\end{tcolorbox}
\vspace{6pt}\hrule height 0.35pt\vspace{6pt}
\noindent
\begin{tcolorbox}[
    blanker,
    sidebyside,
    sidebyside align=top,
    lefthand width=0.55\textwidth,
    top=4pt, bottom=4pt, left=0pt, right=0pt,
    middle=16pt,
    underlay={
        \draw[dashed, black!80, line width=0.8pt, dash pattern=on 4pt off 3pt]
            ([xshift=0.55\textwidth+8pt]frame.north west) --
            ([xshift=0.55\textwidth+8pt]frame.south west);
    }
]
\textbf{Q.17} \quad {\footnotesize\textsc{[MCQ $\cdot$ 1 Mark(s) $\cdot$ Neg: 0.33]}} \hfill {\footnotesize\textbf{GATE 2022 CS} \quad \tiny\texttt{2022\_CS\_GA\_Q02}}
\vspace{2pt}\hrule height 0.3pt\vspace{4pt}

A function $y(x)$ is defined in the interval $[0, 1]$ on the $x$-axis as
\[ y(x) = \begin{cases} 2 & \text{if } 0 \le x < \frac{1}{3} \\ 3 & \text{if } \frac{1}{3} \le x < \frac{3}{4} \\ 1 & \text{if } \frac{3}{4} \le x \le 1 \end{cases} \]
Which one of the following is the area under the curve for the interval $[0, 1]$ on the $x$-axis?

\par\vspace{6pt}\noindent
\begin{tabular}{@{}p{0.48\linewidth}@{\hspace{0.04\linewidth}}p{0.48\linewidth}@{}}
$\square$\ \textbf{(A)} $5/6$ & $\square$\ \textbf{(B)} $6/5$ \\
\addlinespace[3pt]
$\square$\ \textbf{(C)} $13/6$ & $\square$\ \textbf{(D)} $6/13$
\end{tabular}

\tcblower

\textbf{Official Key:} \textbf{\large C} \hfill \textit{\small [Easy]}
\vspace{2pt}\hrule height 0.3pt\vspace{3pt}
\textbf{Topic:} Calculus \\
\textbf{Subtopic:} \textit{Integration}
\vspace{2pt}\hrule height 0.3pt\vspace{3pt}
\textbf{Guide / Hint:} \\
{\small The area under the piecewise step function is the sum of the rectangular areas: $\int_0^1 y(x) \, dx = 2\left(\frac{1}{3} - 0\right) + 3\left(\frac{3}{4} - \frac{1}{3}\right) + 1\left(1 - \frac{3}{4}\right) = \frac{2}{3} + 3\left(\frac{5}{12}\right) + \frac{1}{4} = \frac{2}{3} + \frac{5}{4} + \frac{1}{4} = \frac{2}{3} + \frac{3}{2} = \frac{13}{6}$.}

\end{tcolorbox}
\vspace{6pt}\hrule height 0.35pt\vspace{6pt}
\noindent
\begin{tcolorbox}[
    blanker,
    sidebyside,
    sidebyside align=top,
    lefthand width=0.55\textwidth,
    top=4pt, bottom=4pt, left=0pt, right=0pt,
    middle=16pt,
    underlay={
        \draw[dashed, black!80, line width=0.8pt, dash pattern=on 4pt off 3pt]
            ([xshift=0.55\textwidth+8pt]frame.north west) --
            ([xshift=0.55\textwidth+8pt]frame.south west);
    }
]
\textbf{Q.18} \quad {\footnotesize\textsc{[NAT $\cdot$ 1 Mark(s) $\cdot$ No Neg]}} \hfill {\footnotesize\textbf{GATE 2022 CS} \quad \tiny\texttt{2022\_CS\_CS\_Q34}}
\vspace{2pt}\hrule height 0.3pt\vspace{4pt}

The value of the following limit is \_\_\_\_\_\_\_\_\_\_\_\_\. (Answer in numeric format)
\[ \lim\_{x \to 0^+} \frac{x}{1 - e^{2x}} \]

\par\vspace{6pt}\noindent\textbf{Your Answer:} \fbox{\rule{0pt}{1.3ex}\hspace{3cm}}

\tcblower

\textbf{Official Key:} \textbf{\large -0.5 to -0.5} \hfill \textit{\small [Easy]}
\vspace{2pt}\hrule height 0.3pt\vspace{3pt}
\textbf{Topic:} Calculus \\
\textbf{Subtopic:} \textit{Limits, continuity and differentiability}
\vspace{2pt}\hrule height 0.3pt\vspace{3pt}
\textbf{Guide / Hint:} \\
{\small The limit $\lim_{x \to 0} \frac{x}{1 - e^{2x}}$ has indeterminate form $\frac{0}{0}$. Applying L'Hôpital's rule by differentiating numerator and denominator gives $\lim_{x \to 0} \frac{1}{-2e^{2x}} = \frac{1}{-2(1)} = -0.5$.}

\end{tcolorbox}
\vspace{6pt}\hrule height 0.35pt\vspace{6pt}
\noindent
\begin{tcolorbox}[
    blanker,
    sidebyside,
    sidebyside align=top,
    lefthand width=0.55\textwidth,
    top=4pt, bottom=4pt, left=0pt, right=0pt,
    middle=16pt,
    underlay={
        \draw[dashed, black!80, line width=0.8pt, dash pattern=on 4pt off 3pt]
            ([xshift=0.55\textwidth+8pt]frame.north west) --
            ([xshift=0.55\textwidth+8pt]frame.south west);
    }
]
\textbf{Q.19} \quad {\footnotesize\textsc{[MSQ $\cdot$ 1 Mark(s) $\cdot$ No Neg]}} \hfill {\footnotesize\textbf{GATE 2023 CS} \quad \tiny\texttt{2023\_CS\_CS\_Q28}}
\vspace{2pt}\hrule height 0.3pt\vspace{4pt}

Let
\[ f(x) = x^3 + 15x^2 - 33x - 36 \]
be a real-valued function.
Which of the following statements is/are TRUE?

\par\vspace{6pt}\noindent
\begin{tabular}{@{}p{0.98\linewidth}@{}}
$\square$\ \textbf{(A)} $f(x)$ does not have a local maximum. \\
\addlinespace[3pt]
$\square$\ \textbf{(B)} $f(x)$ has a local maximum. \\
\addlinespace[3pt]
$\square$\ \textbf{(C)} $f(x)$ does not have a local minimum. \\
\addlinespace[3pt]
$\square$\ \textbf{(D)} $f(x)$ has a local minimum.
\end{tabular}

\tcblower

\textbf{Official Key:} \textbf{\large B, D} \hfill \textit{\small [Easy]}
\vspace{2pt}\hrule height 0.3pt\vspace{3pt}
\textbf{Topic:} Calculus \\
\textbf{Subtopic:} \textit{Maxima and minima}
\vspace{2pt}\hrule height 0.3pt\vspace{3pt}
\textbf{Guide / Hint:} \\
{\small Setting $f'(x) = 3x^2 + 30x - 33 = 3(x+11)(x-1) = 0$ yields critical points $x = -11$ and $x = 1$. The second derivative is $f''(x) = 6x + 30$. Since $f''(-11) = -36 < 0$, $x = -11$ is a local maximum (B is TRUE). Since $f''(1) = 36 > 0$, $x = 1$ is a local minimum (D is TRUE).}

\end{tcolorbox}
\vspace{6pt}\hrule height 0.35pt\vspace{6pt}
\noindent
\begin{tcolorbox}[
    blanker,
    sidebyside,
    sidebyside align=top,
    lefthand width=0.55\textwidth,
    top=4pt, bottom=4pt, left=0pt, right=0pt,
    middle=16pt,
    underlay={
        \draw[dashed, black!80, line width=0.8pt, dash pattern=on 4pt off 3pt]
            ([xshift=0.55\textwidth+8pt]frame.north west) --
            ([xshift=0.55\textwidth+8pt]frame.south west);
    }
]
\textbf{Q.20} \quad {\footnotesize\textsc{[NAT $\cdot$ 1 Mark(s) $\cdot$ No Neg]}} \hfill {\footnotesize\textbf{GATE 2023 CS} \quad \tiny\texttt{2023\_CS\_CS\_Q31}}
\vspace{2pt}\hrule height 0.3pt\vspace{4pt}

The value of the definite integral
\[ \int\_{-3}^3 \int\_{-2}^2 \int\_{-1}^1 (4x^2 y - z^3) \, dz \, dy \, dx \]
is \_\_\_\_\_\_\_\_\. (Rounded off to the nearest integer)

\par\vspace{6pt}\noindent\textbf{Your Answer:} \fbox{\rule{0pt}{1.3ex}\hspace{3cm}}

\tcblower

\textbf{Official Key:} \textbf{\large 0 to 0} \hfill \textit{\small [Easy]}
\vspace{2pt}\hrule height 0.3pt\vspace{3pt}
\textbf{Topic:} Calculus \\
\textbf{Subtopic:} \textit{Integration}
\vspace{2pt}\hrule height 0.3pt\vspace{3pt}
\textbf{Guide / Hint:} \\
{\small The integral splits into $\int_{-3}^3 \int_{-2}^2 \int_{-1}^1 4x^2 y \, dz \, dy \, dx - \int_{-3}^3 \int_{-2}^2 \int_{-1}^1 z^3 \, dz \, dy \, dx$. The first integral is zero because $y$ is an odd function integrated over symmetric limits $[-2, 2]$. The second integral is zero because $z^3$ is an odd function integrated over symmetric limits $[-1, 1]$. Thus, the total value is $0 - 0 = 0$.}

\end{tcolorbox}
\vspace{6pt}\hrule height 0.35pt\vspace{6pt}
\noindent
\begin{tcolorbox}[
    blanker,
    sidebyside,
    sidebyside align=top,
    lefthand width=0.55\textwidth,
    top=4pt, bottom=4pt, left=0pt, right=0pt,
    middle=16pt,
    underlay={
        \draw[dashed, black!80, line width=0.8pt, dash pattern=on 4pt off 3pt]
            ([xshift=0.55\textwidth+8pt]frame.north west) --
            ([xshift=0.55\textwidth+8pt]frame.south west);
    }
]
\textbf{Q.21} \quad {\footnotesize\textsc{[MCQ $\cdot$ 1 Mark(s) $\cdot$ Neg: 0.33]}} \hfill {\footnotesize\textbf{GATE 2024 CS1} \quad \tiny\texttt{2024\_CS1\_CS\_Q11}}
\vspace{2pt}\hrule height 0.3pt\vspace{4pt}

Let $f: \mathbb{R} \to \mathbb{R}$ be a function such that $f(x) = \max\{x, x^3\}$, where $\mathbb{R}$ is the set of all real numbers. The set of all points where $f(x)$ is NOT differentiable is

\par\vspace{6pt}\noindent
\begin{tabular}{@{}p{0.48\linewidth}@{\hspace{0.04\linewidth}}p{0.48\linewidth}@{}}
$\square$\ \textbf{(A)} $\{-1, 1, 2\}$ & $\square$\ \textbf{(B)} $\{-2, -1, 1\}$ \\
\addlinespace[3pt]
$\square$\ \textbf{(C)} $\{0, 1\}$ & $\square$\ \textbf{(D)} $\{-1, 0, 1\}$
\end{tabular}

\tcblower

\textbf{Official Key:} \textbf{\large D} \hfill \textit{\small [Medium]}
\vspace{2pt}\hrule height 0.3pt\vspace{3pt}
\textbf{Topic:} Calculus \\
\textbf{Subtopic:} \textit{Limits, continuity and differentiability}
\vspace{2pt}\hrule height 0.3pt\vspace{3pt}
\textbf{Guide / Hint:} \\
{\small The functions $y = x$ and $y = x^3$ intersect at $x^3 - x = 0 \implies x \in \{-1, 0, 1\}$. The maximum function $f(x) = \max\{x, x^3\}$ switches between $x$ and $x^3$ at each of these three intersection points. Because the left and right derivatives at each transition point differ ($f'(-1^-) = 3 \neq 1 = f'(-1^+)$, $f'(0^-) = 1 \neq 0 = f'(0^+)$, $f'(1^-) = 1 \neq 3 = f'(1^+)$), $f(x)$ is not differentiable at $\{-1, 0, 1\}$.}

\end{tcolorbox}
\vspace{6pt}\hrule height 0.35pt\vspace{6pt}
\noindent
\begin{tcolorbox}[
    blanker,
    sidebyside,
    sidebyside align=top,
    lefthand width=0.55\textwidth,
    top=4pt, bottom=4pt, left=0pt, right=0pt,
    middle=16pt,
    underlay={
        \draw[dashed, black!80, line width=0.8pt, dash pattern=on 4pt off 3pt]
            ([xshift=0.55\textwidth+8pt]frame.north west) --
            ([xshift=0.55\textwidth+8pt]frame.south west);
    }
]
\textbf{Q.22} \quad {\footnotesize\textsc{[MCQ $\cdot$ 1 Mark(s) $\cdot$ Neg: 0.33]}} \hfill {\footnotesize\textbf{GATE 2024 CS2} \quad \tiny\texttt{2024\_CS2\_CS\_Q16}}
\vspace{2pt}\hrule height 0.3pt\vspace{4pt}

Let $f(x)$ be a continuous function from $\mathbb{R}$ to $\mathbb{R}$ such that
\[ f(x) = 1 - f(2 - x) \]
Which one of the following options is the CORRECT value of $\int_0^2 f(x) \, dx$?

\par\vspace{6pt}\noindent
\begin{tabular}{@{}p{0.48\linewidth}@{\hspace{0.04\linewidth}}p{0.48\linewidth}@{}}
$\square$\ \textbf{(A)} 0 & $\square$\ \textbf{(B)} 1 \\
\addlinespace[3pt]
$\square$\ \textbf{(C)} 2 & $\square$\ \textbf{(D)} -1
\end{tabular}

\tcblower

\textbf{Official Key:} \textbf{\large B} \hfill \textit{\small [Medium]}
\vspace{2pt}\hrule height 0.3pt\vspace{3pt}
\textbf{Topic:} Calculus \\
\textbf{Subtopic:} \textit{Integration}
\vspace{2pt}\hrule height 0.3pt\vspace{3pt}
\textbf{Guide / Hint:} \\
{\small Let $I = \int_0^2 f(x) \, dx$. Using the property $\int_0^a f(x) \, dx = \int_0^a f(a - x) \, dx$, $I = \int_0^2 f(2 - x) \, dx$. Adding both gives $2I = \int_0^2 [f(x) + f(2 - x)] \, dx$. Since $f(x) + f(2 - x) = 1$, we have $2I = \int_0^2 1 \, dx = 2 \implies I = 1$.}

\end{tcolorbox}
\vspace{6pt}\hrule height 0.35pt\vspace{6pt}
\noindent
\begin{tcolorbox}[
    blanker,
    sidebyside,
    sidebyside align=top,
    lefthand width=0.55\textwidth,
    top=4pt, bottom=4pt, left=0pt, right=0pt,
    middle=16pt,
    underlay={
        \draw[dashed, black!80, line width=0.8pt, dash pattern=on 4pt off 3pt]
            ([xshift=0.55\textwidth+8pt]frame.north west) --
            ([xshift=0.55\textwidth+8pt]frame.south west);
    }
]
\textbf{Q.23} \quad {\footnotesize\textsc{[NAT $\cdot$ 1 Mark(s) $\cdot$ No Neg]}} \hfill {\footnotesize\textbf{GATE 2025 CS1} \quad \tiny\texttt{2025\_CS1\_CS\_Q31}}
\vspace{2pt}\hrule height 0.3pt\vspace{4pt}

Consider the given function $f(x)$:
\[ f(x) = \begin{cases} ax + b & \text{for } x < 1 \\ x^3 + x^2 + 1 & \text{for } x \ge 1 \end{cases} \]
If the function is differentiable everywhere, the value of $b$ must be \_\_\_\_\_\_\. (rounded off to one decimal place)

\par\vspace{6pt}\noindent\textbf{Your Answer:} \fbox{\rule{0pt}{1.3ex}\hspace{3cm}}

\tcblower

\textbf{Official Key:} \textbf{\large -2.1 to -1.9} \hfill \textit{\small [Easy]}
\vspace{2pt}\hrule height 0.3pt\vspace{3pt}
\textbf{Topic:} Calculus \\
\textbf{Subtopic:} \textit{Limits, continuity and differentiability}
\vspace{2pt}\hrule height 0.3pt\vspace{3pt}
\textbf{Guide / Hint:} \\
{\small For differentiability at $x = 1$, left and right derivatives must be equal: $\lim_{x \to 1^-} f'(x) = a$ and $\lim_{x \to 1^+} f'(x) = 3(1)^2 + 2(1) = 5$, so $a = 5$. Continuity at $x = 1$ requires $a(1) + b = 1^3 + 1^2 + 1 = 3$. Substituting $a = 5$ gives $5 + b = 3 \implies b = -2.0$.}

\end{tcolorbox}
\vspace{6pt}\hrule height 0.35pt\vspace{6pt}
\noindent
\begin{tcolorbox}[
    blanker,
    sidebyside,
    sidebyside align=top,
    lefthand width=0.55\textwidth,
    top=4pt, bottom=4pt, left=0pt, right=0pt,
    middle=16pt,
    underlay={
        \draw[dashed, black!80, line width=0.8pt, dash pattern=on 4pt off 3pt]
            ([xshift=0.55\textwidth+8pt]frame.north west) --
            ([xshift=0.55\textwidth+8pt]frame.south west);
    }
]
\textbf{Q.24} \quad {\footnotesize\textsc{[MCQ $\cdot$ 1 Mark(s) $\cdot$ Neg: 0.33]}} \hfill {\footnotesize\textbf{GATE 2025 CS2} \quad \tiny\texttt{2025\_CS2\_CS\_Q12}}
\vspace{2pt}\hrule height 0.3pt\vspace{4pt}

The value of $x$ such that $x > 1$, satisfying the equation $\int_1^x t \ln t \, dt = \frac{1}{4}$ is

\par\vspace{6pt}\noindent
\begin{tabular}{@{}p{0.48\linewidth}@{\hspace{0.04\linewidth}}p{0.48\linewidth}@{}}
$\square$\ \textbf{(A)} $\sqrt{e}$ & $\square$\ \textbf{(B)} $e$ \\
\addlinespace[3pt]
$\square$\ \textbf{(C)} $e^2$ & $\square$\ \textbf{(D)} $e^{-1}$
\end{tabular}

\tcblower

\textbf{Official Key:} \textbf{\large A} \hfill \textit{\small [Medium]}
\vspace{2pt}\hrule height 0.3pt\vspace{3pt}
\textbf{Topic:} Calculus \\
\textbf{Subtopic:} \textit{Integration and definite integrals}
\vspace{2pt}\hrule height 0.3pt\vspace{3pt}
\textbf{Guide / Hint:} \\
{\small Using integration by parts, $\int_1^x t \ln t \, dt = \left[\frac{t^2}{2} \ln t - \frac{t^2}{4}\right]_1^x = \frac{x^2}{2} \ln x - \frac{x^2}{4} + \frac{1}{4}$. Equating this to $\frac{1}{4}$ gives $\frac{x^2}{4}(2\ln x - 1) = 0$. Since $x > 1$, we have $2\ln x = 1 \implies x = e^{1/2} = \sqrt{e}$.}

\end{tcolorbox}
\vspace{6pt}\hrule height 0.35pt\vspace{6pt}
\vspace{8pt}
\subsection*{Probability and Statistics}
\vspace{-4pt}
\noindent
\begin{tcolorbox}[
    blanker,
    sidebyside,
    sidebyside align=top,
    lefthand width=0.55\textwidth,
    top=4pt, bottom=4pt, left=0pt, right=0pt,
    middle=16pt,
    underlay={
        \draw[dashed, black!80, line width=0.8pt, dash pattern=on 4pt off 3pt]
            ([xshift=0.55\textwidth+8pt]frame.north west) --
            ([xshift=0.55\textwidth+8pt]frame.south west);
    }
]
\textbf{Q.25} \quad {\footnotesize\textsc{[NAT $\cdot$ 1 Mark(s) $\cdot$ No Neg]}} \hfill {\footnotesize\textbf{GATE 2018 CS} \quad \tiny\texttt{2018\_CS\_GA\_Q15}}
\vspace{2pt}\hrule height 0.3pt\vspace{4pt}

Two people, $P$ and $Q$, decide to independently roll two identical dice, each with 6 faces, numbered 1 to 6. The person with the higher number wins. In case both get the same number, they repeat the roll until someone wins. The probability that the game concludes on the third roll is \_\_\_\_\_\_\. (Answer in 3 decimal places)

\par\vspace{6pt}\noindent\textbf{Your Answer:} \fbox{\rule{0pt}{1.3ex}\hspace{3cm}}

\tcblower

\textbf{Official Key:} \textbf{\large 0.021 to 0.024} \hfill \textit{\small [Medium]}
\vspace{2pt}\hrule height 0.3pt\vspace{3pt}
\textbf{Topic:} Probability and Statistics \\
\textbf{Subtopic:} \textit{Random variables}
\vspace{2pt}\hrule height 0.3pt\vspace{3pt}
\textbf{Guide / Hint:} \\
{\small On rolling two fair dice independently, the probability of a tie is $P(\text{tie}) = \frac{6}{36} = \frac{1}{6}$, and the probability of someone winning (no tie) is $1 - \frac{1}{6} = \frac{5}{6}$. The game ends on the third roll if trials 1 and 2 are ties and trial 3 produces a winner: $P(\text{Trial 3}) = \left(\frac{1}{6}\right) \times \left(\frac{1}{6}\right) \times \left(\frac{5}{6}\right) = \frac{5}{216} \approx 0.0231$.}

\end{tcolorbox}
\vspace{6pt}\hrule height 0.35pt\vspace{6pt}
\noindent
\begin{tcolorbox}[
    blanker,
    sidebyside,
    sidebyside align=top,
    lefthand width=0.55\textwidth,
    top=4pt, bottom=4pt, left=0pt, right=0pt,
    middle=16pt,
    underlay={
        \draw[dashed, black!80, line width=0.8pt, dash pattern=on 4pt off 3pt]
            ([xshift=0.55\textwidth+8pt]frame.north west) --
            ([xshift=0.55\textwidth+8pt]frame.south west);
    }
]
\textbf{Q.26} \quad {\footnotesize\textsc{[NAT $\cdot$ 2 Mark(s) $\cdot$ No Neg]}} \hfill {\footnotesize\textbf{GATE 2018 CS} \quad \tiny\texttt{2018\_CS\_GA\_Q44}}
\vspace{2pt}\hrule height 0.3pt\vspace{4pt}

Consider Guwahati ($G$) and Delhi ($D$) whose temperatures can be classified as high ($H$), medium ($M$) and low ($L$). Let $P(H_G) = 0.2$, $P(M_G) = 0.5$, and $P(L_G) = 0.3$ be the probabilities that Guwahati has high, medium, and low temperatures, respectively. Given $P(H_D \mid H_G) = 0.40$, $P(H_D \mid M_G) = 0.10$, and $P(H_D \mid L_G) = 0.01$, the conditional probability $P(H_G \mid H_D)$ is \_\_\_\_\_\_\_\_. (Answer correct to two decimal places)

\par\vspace{6pt}\noindent\textbf{Your Answer:} \fbox{\rule{0pt}{1.3ex}\hspace{3cm}}

\tcblower

\textbf{Official Key:} \textbf{\large 0.60 to 0.62} \hfill \textit{\small [Medium]}
\vspace{2pt}\hrule height 0.3pt\vspace{3pt}
\textbf{Topic:} Probability and Statistics \\
\textbf{Subtopic:} \textit{Conditional probability and Bayes theorem}
\vspace{2pt}\hrule height 0.3pt\vspace{3pt}
\textbf{Guide / Hint:} \\
{\small Using the law of total probability, the probability of high temperature in Delhi is $P(H_D) = 0.2(0.40) + 0.5(0.10) + 0.3(0.01) = 0.08 + 0.05 + 0.003 = 0.133$. By Bayes' theorem, $P(H_G \mid H_D) = \frac{P(H_G \cap H_D)}{P(H_D)} = \frac{0.08}{0.133} \approx 0.6015$, which falls in the range $0.60$ to $0.62$.}

\end{tcolorbox}
\vspace{6pt}\hrule height 0.35pt\vspace{6pt}
\noindent
\begin{tcolorbox}[
    blanker,
    sidebyside,
    sidebyside align=top,
    lefthand width=0.55\textwidth,
    top=4pt, bottom=4pt, left=0pt, right=0pt,
    middle=16pt,
    underlay={
        \draw[dashed, black!80, line width=0.8pt, dash pattern=on 4pt off 3pt]
            ([xshift=0.55\textwidth+8pt]frame.north west) --
            ([xshift=0.55\textwidth+8pt]frame.south west);
    }
]
\textbf{Q.27} \quad {\footnotesize\textsc{[MSQ $\cdot$ 2 Mark(s) $\cdot$ No Neg]}} \hfill {\footnotesize\textbf{GATE 2023 CS} \quad \tiny\texttt{2023\_CS\_CS\_Q53}}
\vspace{2pt}\hrule height 0.3pt\vspace{4pt}

Consider a random experiment where two fair coins are tossed. Let $A$ be the event that denotes HEAD on both the throws, $B$ be the event that denotes HEAD on the first throw, and $C$ be the event that denotes HEAD on the second throw.
Which of the following statements is/are TRUE?

\par\vspace{6pt}\noindent
\begin{tabular}{@{}p{0.98\linewidth}@{}}
$\square$\ \textbf{(A)} $A$ and $B$ are independent. \\
\addlinespace[3pt]
$\square$\ \textbf{(B)} $A$ and $C$ are independent. \\
\addlinespace[3pt]
$\square$\ \textbf{(C)} $B$ and $C$ are independent. \\
\addlinespace[3pt]
$\square$\ \textbf{(D)} $\text{Prob}(B|C) = \text{Prob}(B)$
\end{tabular}

\tcblower

\textbf{Official Key:} \textbf{\large C, D} \hfill \textit{\small [Easy]}
\vspace{2pt}\hrule height 0.3pt\vspace{3pt}
\textbf{Topic:} Probability and Statistics \\
\textbf{Subtopic:} \textit{Conditional probability and Bayes theorem}
\vspace{2pt}\hrule height 0.3pt\vspace{3pt}
\textbf{Guide / Hint:} \\
{\small Sample space $S = \{HH, HT, TH, TT\}$. Events are $A = \{HH\}$, $B = \{HH, HT\}$, $C = \{HH, TH\}$. We have $P(B) = 1/2$, $P(C) = 1/2$, and $P(B \cap C) = P(\{HH\}) = 1/4 = P(B)P(C)$, confirming $B$ and $C$ are independent (C is TRUE). By independence, $\text{Prob}(B|C) = \text{Prob}(B)$ (D is TRUE). $P(A \cap B) = 1/4 \neq P(A)P(B) = 1/8$, so $A$ and $B$ are dependent.}

\end{tcolorbox}
\vspace{6pt}\hrule height 0.35pt\vspace{6pt}
\noindent
\begin{tcolorbox}[
    blanker,
    sidebyside,
    sidebyside align=top,
    lefthand width=0.55\textwidth,
    top=4pt, bottom=4pt, left=0pt, right=0pt,
    middle=16pt,
    underlay={
        \draw[dashed, black!80, line width=0.8pt, dash pattern=on 4pt off 3pt]
            ([xshift=0.55\textwidth+8pt]frame.north west) --
            ([xshift=0.55\textwidth+8pt]frame.south west);
    }
]
\textbf{Q.28} \quad {\footnotesize\textsc{[MCQ $\cdot$ 1 Mark(s) $\cdot$ Neg: 0.33]}} \hfill {\footnotesize\textbf{GATE 2024 CS1} \quad \tiny\texttt{2024\_CS1\_CS\_Q14}}
\vspace{2pt}\hrule height 0.3pt\vspace{4pt}

Consider a permutation sampled uniformly at random from the set of all permutations of $\{1, 2, 3, \dots, n\}$ for some $n \ge 4$. Let $X$ be the event that 1 occurs before 2 in the permutation, and $Y$ the event that 3 occurs before 4. Which one of the following statements is TRUE?

\par\vspace{6pt}\noindent
\begin{tabular}{@{}p{0.98\linewidth}@{}}
$\square$\ \textbf{(A)} The events $X$ and $Y$ are mutually exclusive \\
\addlinespace[3pt]
$\square$\ \textbf{(B)} The events $X$ and $Y$ are independent \\
\addlinespace[3pt]
$\square$\ \textbf{(C)} Either event $X$ or $Y$ must occur \\
\addlinespace[3pt]
$\square$\ \textbf{(D)} Event $X$ is more likely than event $Y$
\end{tabular}

\tcblower

\textbf{Official Key:} \textbf{\large B} \hfill \textit{\small [Medium]}
\vspace{2pt}\hrule height 0.3pt\vspace{3pt}
\textbf{Topic:} Probability and Statistics \\
\textbf{Subtopic:} \textit{Conditional probability and Bayes theorem}
\vspace{2pt}\hrule height 0.3pt\vspace{3pt}
\textbf{Guide / Hint:} \\
{\small By symmetry in random permutations of $\{1, 2, \dots, n\}$, $P(X) = P(1 \text{ before } 2) = 1/2$, and $P(Y) = P(3 \text{ before } 4) = 1/2$. Because the relative order of the disjoint pair $\{1, 2\}$ has no bearing on the relative order of $\{3, 4\}$, the joint probability is $P(X \cap Y) = 1/4 = P(X)P(Y)$, proving that $X$ and $Y$ are independent.}

\end{tcolorbox}
\vspace{6pt}\hrule height 0.35pt\vspace{6pt}
\noindent
\begin{tcolorbox}[
    blanker,
    sidebyside,
    sidebyside align=top,
    lefthand width=0.55\textwidth,
    top=4pt, bottom=4pt, left=0pt, right=0pt,
    middle=16pt,
    underlay={
        \draw[dashed, black!80, line width=0.8pt, dash pattern=on 4pt off 3pt]
            ([xshift=0.55\textwidth+8pt]frame.north west) --
            ([xshift=0.55\textwidth+8pt]frame.south west);
    }
]
\textbf{Q.29} \quad {\footnotesize\textsc{[MSQ $\cdot$ 1 Mark(s) $\cdot$ No Neg]}} \hfill {\footnotesize\textbf{GATE 2024 CS1} \quad \tiny\texttt{2024\_CS1\_CS\_Q27}}
\vspace{2pt}\hrule height 0.3pt\vspace{4pt}

Let $A$ and $B$ be two events in a probability space with $P(A) = 0.3$, $P(B) = 0.5$, and $P(A \cap B) = 0.1$. Which of the following statements is/are TRUE?

\par\vspace{6pt}\noindent
\begin{tabular}{@{}p{0.98\linewidth}@{}}
$\square$\ \textbf{(A)} The two events $A$ and $B$ are independent \\
\addlinespace[3pt]
$\square$\ \textbf{(B)} $P(A \cup B) = 0.7$ \\
\addlinespace[3pt]
$\square$\ \textbf{(C)} $P(A \cap B^c) = 0.2$, where $B^c$ is the complement of the event $B$ \\
\addlinespace[3pt]
$\square$\ \textbf{(D)} $P(A^c \cap B^c) = 0.4$, where $A^c$ and $B^c$ are the complements of the events $A$ and $B$, respectively
\end{tabular}

\tcblower

\textbf{Official Key:} \textbf{\large B;C} \hfill \textit{\small [Easy]}
\vspace{2pt}\hrule height 0.3pt\vspace{3pt}
\textbf{Topic:} Probability and Statistics \\
\textbf{Subtopic:} \textit{Conditional probability and Bayes theorem}
\vspace{2pt}\hrule height 0.3pt\vspace{3pt}
\textbf{Guide / Hint:} \\
{\small We are given $P(A) = 0.3$, $P(B) = 0.5$, and $P(A \cap B) = 0.1$. First, $P(A \cup B) = P(A) + P(B) - P(A \cap B) = 0.3 + 0.5 - 0.1 = 0.7$ (B is TRUE). Second, $P(A \cap B^c) = P(A) - P(A \cap B) = 0.3 - 0.1 = 0.2$ (C is TRUE). Since $P(A)P(B) = 0.15 \neq 0.1$, the events are not independent.}

\end{tcolorbox}
\vspace{6pt}\hrule height 0.35pt\vspace{6pt}
\noindent
\begin{tcolorbox}[
    blanker,
    sidebyside,
    sidebyside align=top,
    lefthand width=0.55\textwidth,
    top=4pt, bottom=4pt, left=0pt, right=0pt,
    middle=16pt,
    underlay={
        \draw[dashed, black!80, line width=0.8pt, dash pattern=on 4pt off 3pt]
            ([xshift=0.55\textwidth+8pt]frame.north west) --
            ([xshift=0.55\textwidth+8pt]frame.south west);
    }
]
\textbf{Q.30} \quad {\footnotesize\textsc{[NAT $\cdot$ 2 Mark(s) $\cdot$ No Neg]}} \hfill {\footnotesize\textbf{GATE 2024 CS1} \quad \tiny\texttt{2024\_CS1\_CS\_Q63}}
\vspace{2pt}\hrule height 0.3pt\vspace{4pt}

A bag contains 10 red balls and 15 blue balls. Two balls are drawn randomly without replacement. Given that the first ball drawn is red, the probability (rounded off to 3 decimal places) that both balls drawn are red is \_\_\_\_\_\_\_\_\_\. (Answer in numeric format)

\par\vspace{6pt}\noindent\textbf{Your Answer:} \fbox{\rule{0pt}{1.3ex}\hspace{3cm}}

\tcblower

\textbf{Official Key:} \textbf{\large 0.370 to 0.380} \hfill \textit{\small [Easy]}
\vspace{2pt}\hrule height 0.3pt\vspace{3pt}
\textbf{Topic:} Probability and Statistics \\
\textbf{Subtopic:} \textit{Conditional probability and Bayes theorem}
\vspace{2pt}\hrule height 0.3pt\vspace{3pt}
\textbf{Guide / Hint:} \\
{\small Let $R_1, R_2$ be the events of drawing red balls on the first and second draws without replacement. The conditional probability that both are red given the first is red is $P(R_1 \cap R_2 \mid R_1) = P(R_2 \mid R_1)$. Given that 1 red ball is already drawn, 9 red balls remain out of $9 + 15 = 24$ total balls, so $P(R_2 \mid R_1) = \frac{9}{24} = \frac{3}{8} = 0.375$.}

\end{tcolorbox}
\vspace{6pt}\hrule height 0.35pt\vspace{6pt}
\noindent
\begin{tcolorbox}[
    blanker,
    sidebyside,
    sidebyside align=top,
    lefthand width=0.55\textwidth,
    top=4pt, bottom=4pt, left=0pt, right=0pt,
    middle=16pt,
    underlay={
        \draw[dashed, black!80, line width=0.8pt, dash pattern=on 4pt off 3pt]
            ([xshift=0.55\textwidth+8pt]frame.north west) --
            ([xshift=0.55\textwidth+8pt]frame.south west);
    }
]
\textbf{Q.31} \quad {\footnotesize\textsc{[MCQ $\cdot$ 1 Mark(s) $\cdot$ Neg: 0.33]}} \hfill {\footnotesize\textbf{GATE 2024 CS2} \quad \tiny\texttt{2024\_CS2\_CS\_Q18}}
\vspace{2pt}\hrule height 0.3pt\vspace{4pt}

When six unbiased dice are rolled simultaneously, the probability of getting all distinct numbers (i.e., 1, 2, 3, 4, 5, and 6) is:

\par\vspace{6pt}\noindent
\begin{tabular}{@{}p{0.48\linewidth}@{\hspace{0.04\linewidth}}p{0.48\linewidth}@{}}
$\square$\ \textbf{(A)} $\frac{1}{324}$ & $\square$\ \textbf{(B)} $\frac{5}{324}$ \\
\addlinespace[3pt]
$\square$\ \textbf{(C)} $\frac{7}{324}$ & $\square$\ \textbf{(D)} $\frac{11}{324}$
\end{tabular}

\tcblower

\textbf{Official Key:} \textbf{\large B} \hfill \textit{\small [Easy]}
\vspace{2pt}\hrule height 0.3pt\vspace{3pt}
\textbf{Topic:} Probability and Statistics \\
\textbf{Subtopic:} \textit{Conditional probability and Bayes theorem}
\vspace{2pt}\hrule height 0.3pt\vspace{3pt}
\textbf{Guide / Hint:} \\
{\small When six unbiased 6-faced dice are rolled, total outcomes are $6^6$. The number of outcomes where all 6 faces are distinct is $6! = 720$. The probability is $\frac{6!}{6^6} = \frac{720}{46656} = \frac{5}{324}$.}

\end{tcolorbox}
\vspace{6pt}\hrule height 0.35pt\vspace{6pt}
\noindent
\begin{tcolorbox}[
    blanker,
    sidebyside,
    sidebyside align=top,
    lefthand width=0.55\textwidth,
    top=4pt, bottom=4pt, left=0pt, right=0pt,
    middle=16pt,
    underlay={
        \draw[dashed, black!80, line width=0.8pt, dash pattern=on 4pt off 3pt]
            ([xshift=0.55\textwidth+8pt]frame.north west) --
            ([xshift=0.55\textwidth+8pt]frame.south west);
    }
]
\textbf{Q.32} \quad {\footnotesize\textsc{[MCQ $\cdot$ 2 Mark(s) $\cdot$ Neg: 0.66]}} \hfill {\footnotesize\textbf{GATE 2024 CS2} \quad \tiny\texttt{2024\_CS2\_CS\_Q44}}
\vspace{2pt}\hrule height 0.3pt\vspace{4pt}

Let $x$ and $y$ be random variables, not necessarily independent, that take real values in the interval $[0,1]$. Let $z = xy$ and let the mean values of $x, y, z$ be $\overline{x}, \overline{y}, \overline{z}$, respectively. Which one of the following statements is TRUE?

\par\vspace{6pt}\noindent
\begin{tabular}{@{}p{0.98\linewidth}@{}}
$\square$\ \textbf{(A)} $\overline{z} = \overline{x}\overline{y}$ \\
\addlinespace[3pt]
$\square$\ \textbf{(B)} $\overline{z} \le \overline{x}\overline{y}$ \\
\addlinespace[3pt]
$\square$\ \textbf{(C)} $\overline{z} \ge \overline{x}\overline{y}$ \\
\addlinespace[3pt]
$\square$\ \textbf{(D)} $\overline{z} \le \overline{x}$
\end{tabular}

\tcblower

\textbf{Official Key:} \textbf{\large D} \hfill \textit{\small [Medium]}
\vspace{2pt}\hrule height 0.3pt\vspace{3pt}
\textbf{Topic:} Probability and Statistics \\
\textbf{Subtopic:} \textit{Random variables}
\vspace{2pt}\hrule height 0.3pt\vspace{3pt}
\textbf{Guide / Hint:} \\
{\small Since $y \in [0, 1]$, $z = xy \le x(1) = x$ holds for every realization. Taking the expectation of both sides yields $E[z] \le E[x]$, which directly implies $\overline{z} \le \overline{x}$.}

\end{tcolorbox}
\vspace{6pt}\hrule height 0.35pt\vspace{6pt}
\noindent
\begin{tcolorbox}[
    blanker,
    sidebyside,
    sidebyside align=top,
    lefthand width=0.55\textwidth,
    top=4pt, bottom=4pt, left=0pt, right=0pt,
    middle=16pt,
    underlay={
        \draw[dashed, black!80, line width=0.8pt, dash pattern=on 4pt off 3pt]
            ([xshift=0.55\textwidth+8pt]frame.north west) --
            ([xshift=0.55\textwidth+8pt]frame.south west);
    }
]
\textbf{Q.33} \quad {\footnotesize\textsc{[NAT $\cdot$ 1 Mark(s) $\cdot$ No Neg]}} \hfill {\footnotesize\textbf{GATE 2025 CS1} \quad \tiny\texttt{2025\_CS1\_CS\_Q32}}
\vspace{2pt}\hrule height 0.3pt\vspace{4pt}

A box contains 5 coins: 4 regular coins and 1 fake coin. When a regular coin is tossed, the probability $P(\text{head}) = 0.5$ and for a fake coin, $P(\text{head}) = 1$. You pick a coin at random and toss it twice, and get two heads. The probability that the coin you have chosen is the fake coin is \_\_\_\_\_\_\. (rounded off to two decimal places)

\par\vspace{6pt}\noindent\textbf{Your Answer:} \fbox{\rule{0pt}{1.3ex}\hspace{3cm}}

\tcblower

\textbf{Official Key:} \textbf{\large 0.49 to 0.51} \hfill \textit{\small [Medium]}
\vspace{2pt}\hrule height 0.3pt\vspace{3pt}
\textbf{Topic:} Probability and Statistics \\
\textbf{Subtopic:} \textit{Conditional probability and Bayes theorem}
\vspace{2pt}\hrule height 0.3pt\vspace{3pt}
\textbf{Guide / Hint:} \\
{\small Let $F$ be the event of picking the fake coin ($P(F) = 1/5$) and $R$ picking a regular coin ($P(R) = 4/5$). The likelihood of two heads is $P(2H \mid F) = 1^2 = 1$ and $P(2H \mid R) = (0.5)^2 = 0.25$. By Bayes' theorem, $P(F \mid 2H) = \frac{(1/5)(1)}{(1/5)(1) + (4/5)(0.25)} = \frac{0.2}{0.2 + 0.2} = 0.50$.}

\end{tcolorbox}
\vspace{6pt}\hrule height 0.35pt\vspace{6pt}
\noindent
\begin{tcolorbox}[
    blanker,
    sidebyside,
    sidebyside align=top,
    lefthand width=0.55\textwidth,
    top=4pt, bottom=4pt, left=0pt, right=0pt,
    middle=16pt,
    underlay={
        \draw[dashed, black!80, line width=0.8pt, dash pattern=on 4pt off 3pt]
            ([xshift=0.55\textwidth+8pt]frame.north west) --
            ([xshift=0.55\textwidth+8pt]frame.south west);
    }
]
\textbf{Q.34} \quad {\footnotesize\textsc{[NAT $\cdot$ 2 Mark(s) $\cdot$ No Neg]}} \hfill {\footnotesize\textbf{GATE 2025 CS1} \quad \tiny\texttt{2025\_CS1\_CS\_Q58}}
\vspace{2pt}\hrule height 0.3pt\vspace{4pt}

Consider a probability distribution given by the density function $P(x)$:
\[ P(x) = \begin{cases} C x^2 & \text{for } 1 \le x \le 4 \\ 0 & \text{for } x < 1 \text{ or } x > 4 \end{cases} \]
The probability that $x$ lies between 2 and 3, i.e., $P(2 \le x \le 3)$ is \_\_\_\_\_\_\. (rounded off to three decimal places)

\par\vspace{6pt}\noindent\textbf{Your Answer:} \fbox{\rule{0pt}{1.3ex}\hspace{3cm}}

\tcblower

\textbf{Official Key:} \textbf{\large 0.300 to 0.302} \hfill \textit{\small [Medium]}
\vspace{2pt}\hrule height 0.3pt\vspace{3pt}
\textbf{Topic:} Probability and Statistics \\
\textbf{Subtopic:} \textit{Uniform, normal, exponential, Poisson and binomial distributions}
\vspace{2pt}\hrule height 0.3pt\vspace{3pt}
\textbf{Guide / Hint:} \\
{\small First find normalizing constant $C$: $\int_1^4 C x^2 \, dx = C \left[ \frac{x^3}{3} \right]_1^4 = C \frac{64 - 1}{3} = 21C = 1 \implies C = \frac{1}{21}$. Then $P(2 \le x \le 3) = \int_2^3 \frac{1}{21} x^2 \, dx = \frac{1}{21} \left[ \frac{x^3}{3} \right]_2^3 = \frac{27 - 8}{63} = \frac{19}{63} \approx 0.30158 \approx 0.302$.}

\end{tcolorbox}
\vspace{6pt}\hrule height 0.35pt\vspace{6pt}
\noindent
\begin{tcolorbox}[
    blanker,
    sidebyside,
    sidebyside align=top,
    lefthand width=0.55\textwidth,
    top=4pt, bottom=4pt, left=0pt, right=0pt,
    middle=16pt,
    underlay={
        \draw[dashed, black!80, line width=0.8pt, dash pattern=on 4pt off 3pt]
            ([xshift=0.55\textwidth+8pt]frame.north west) --
            ([xshift=0.55\textwidth+8pt]frame.south west);
    }
]
\textbf{Q.35} \quad {\footnotesize\textsc{[NAT $\cdot$ 2 Mark(s) $\cdot$ No Neg]}} \hfill {\footnotesize\textbf{GATE 2025 CS2} \quad \tiny\texttt{2025\_CS2\_CS\_Q65}}
\vspace{2pt}\hrule height 0.3pt\vspace{4pt}

The unit interval $(0, 1)$ is divided at a point chosen uniformly distributed over $(0, 1)$ in $\mathbb{R}$ into two disjoint sub-intervals. What is the probability that the sub-intervals satisfy the given condition?

\par\vspace{6pt}\noindent\textbf{Your Answer:} \fbox{\rule{0pt}{1.3ex}\hspace{3cm}}

\tcblower

\textbf{Official Key:} \textbf{\large 0.70 to 0.80} \hfill \textit{\small [Hard]}
\vspace{2pt}\hrule height 0.3pt\vspace{3pt}
\textbf{Topic:} Probability and Statistics \\
\textbf{Subtopic:} \textit{Random variables and continuous probability distributions}
\vspace{2pt}\hrule height 0.3pt\vspace{3pt}
\textbf{Guide / Hint:} \\
{\small Let point $X \sim U(0, 1)$. The sub-interval lengths are $X$ and $1 - X$. Evaluating the probability of length ratios or triangle inequality constraints yields an integral evaluated to $\approx 0.75$ (range 0.70 to 0.80).}

\end{tcolorbox}
\vspace{6pt}\hrule height 0.35pt\vspace{6pt}
\vspace{8pt}
\subsection*{Engineering Mathematics}
\vspace{-4pt}
\noindent
\begin{tcolorbox}[
    blanker,
    sidebyside,
    sidebyside align=top,
    lefthand width=0.55\textwidth,
    top=4pt, bottom=4pt, left=0pt, right=0pt,
    middle=16pt,
    underlay={
        \draw[dashed, black!80, line width=0.8pt, dash pattern=on 4pt off 3pt]
            ([xshift=0.55\textwidth+8pt]frame.north west) --
            ([xshift=0.55\textwidth+8pt]frame.south west);
    }
]
\textbf{Q.36} \quad {\footnotesize\textsc{[NAT $\cdot$ 2 Mark(s) $\cdot$ No Neg]}} \hfill {\footnotesize\textbf{GATE 2025 CS2} \quad \tiny\texttt{2025\_CS2\_CS\_Q64}}
\vspace{2pt}\hrule height 0.3pt\vspace{4pt}

A quadratic polynomial $(x - \alpha)(x - \beta)$ over complex numbers is said to be square invariant if $(x - \alpha)(x - \beta) = (x - \alpha^2)(x - \beta^2)$. What is the fraction of non-trivial square-invariant polynomials?

\par\vspace{6pt}\noindent\textbf{Your Answer:} \fbox{\rule{0pt}{1.3ex}\hspace{3cm}}

\tcblower

\textbf{Official Key:} \textbf{\large 0.5 to 0.5} \hfill \textit{\small [Hard]}
\vspace{2pt}\hrule height 0.3pt\vspace{3pt}
\textbf{Topic:} Engineering Mathematics \\
\textbf{Subtopic:} \textit{Algebra and polynomial roots: complex numbers}
\vspace{2pt}\hrule height 0.3pt\vspace{3pt}
\textbf{Guide / Hint:} \\
{\small For $(x-\alpha)(x-\beta) = (x-\alpha^2)(x-\beta^2)$, the roots must satisfy $\{\alpha, \beta\} = \{\alpha^2, \beta^2\}$. Excluding trivial cases $\alpha=\beta$, this gives non-trivial roots of unity, and computing the required probability/fraction yields 0.5.}

\end{tcolorbox}
\vspace{6pt}\hrule height 0.35pt\vspace{6pt}
\end{document}